% TOP_PROBLEM: {"id": 20000882, "title": "A root-system reduction for the spherical triangle form", "category_id": 2, "category": {"id": 2, "name": "combinatorics", "display_name": "Combinatorics", "description": "Counting problems, graph theory, discrete structures.", "slug": "combinatorics", "order_index": 2, "created_at": "2026-07-31T15:26:25.671Z"}, "status": "partially_solved", "problem_number": "AIM-COMBINATORICS-0007", "set_id": 14}
% FIRST_POSTED: 2026-10-01
% ENTRY_KIND: statement_only
% UnsolvedMath ID 20000882; source catalogue code AIM-COMBINATORICS-0007
% !TeX program = lualatex
% Definition and sources only; this file is not an LLM solution attempt.
\documentclass[11pt,a4paper]{article}
\usepackage[margin=25mm]{geometry}
\usepackage{fontspec}
\setmainfont{lmroman10-regular.otf}[BoldFont=lmroman10-bold.otf,
 ItalicFont=lmroman10-italic.otf,BoldItalicFont=lmroman10-bolditalic.otf]
\usepackage{amsmath,amssymb,mathtools}
\usepackage{unicode-math}
\setmathfont{latinmodern-math.otf}
\AtBeginDocument{\let\setminus\smallsetminus}
\usepackage{xurl,hyperref}
\hypersetup{colorlinks=true,urlcolor=blue}
\setcounter{secnumdepth}{0}
\setlength{\parindent}{0pt}
\setlength{\parskip}{5pt}
\setlength{\emergencystretch}{3em}
\begin{document}
\section*{A root-system reduction for the spherical triangle form}\label{20000882}
{\small UnsolvedMath ID 20000882 · AIM-COMBINATORICS-0007 · Combinatorics · AIM Workshop Problem Lists}
\subsection{Definitions and mathematical statement}
Let $S^n=\{x\in\mathbb R^{n+1}:\|x\|_2=1\}$, $n\geq1$, with normalized rotation-invariant measure $\sigma_n$.
Choose any fixed ordered triple $x_0,y_0,z_0\in S^n$ satisfying $x_0+y_0+z_0=0$.
Its three pairwise inner products are $-1/2$.
For a measurable real function $f$ with $|f|\leq1$, define
\[
 T(f)=\mathbb E_R f(Rx_0)f(Ry_0)f(Rz_0),
\]
where $R$ is Haar-uniform on the orthogonal group $O(n+1)$.
This definition specifies the invariant probability measure on the constrained triples and is independent of the initially chosen triple.

The proposed inverse theorem asks whether for every $\delta\in(0,1)$ there exist $m(\delta)\in\mathbb N$ and $c(\delta)>0$, independent of $n$, with the following property.
Whenever $|T(f)|>\delta$, there are an orthonormal family $u_1,\ldots,u_s$ with $s\leq\min\{m(\delta),n+1\}$ and a measurable function $g:\mathbb R^s\to[-1,1]$ such that
\[
 \left|\int_{S^n}f(x)\,
 g(\langle x,u_1\rangle,\ldots,\langle x,u_s\rangle)
 \,d\sigma_n(x)\right|\geq c(\delta).
\]
Thus both the number of coordinates and the correlation threshold are uniform in ambient dimension.
The coordinates may be chosen after seeing $f$; they need not be the original coordinate axes.
The bound on $g$ fixes the normalization of ``correlation'' and prevents rescaling a negligible correlation into a meaningful one.
\subsection{Short English statement}
Must a bounded function with a large equilateral-triangle average correlate with a function of only boundedly many suitably chosen coordinates?
\subsection{Sources}
\begin{itemize}
\item[S1] ULAM.AI and the UnsolvedMath Contributors, supplied problems.json, record 20000882 (AIM-COMBINATORICS-0007), AIM Workshop Problem Lists. Catalogue identity and source transcription; CC BY 4.0.
\url{https://huggingface.co/datasets/ulamai/UnsolvedMath}.
\item[S2] American Institute of Mathematics, High-dimensional phenomena in discrete analysis, item 2.4. Original workshop question underlying AIM-COMBINATORICS-0007.
\url{http://aimpl.org/highdimdiscrete/2/}.
\end{itemize}
\end{document}
